\documentclass[11pt]{article}

\usepackage[T1]{fontenc}
\usepackage{lmodern}
\usepackage{microtype}
\usepackage[margin=1in]{geometry}
\usepackage{amsmath,amssymb,amsthm,mathtools,bm}
\usepackage{booktabs}
\usepackage{graphicx}
\usepackage{xcolor}
\usepackage[hidelinks]{hyperref}
\usepackage[authoryear,round]{natbib}

\graphicspath{{figures/}}

\theoremstyle{plain}
\newtheorem{theorem}{Theorem}
\newtheorem{lemma}[theorem]{Lemma}
\newtheorem{proposition}[theorem]{Proposition}
\newtheorem{corollary}[theorem]{Corollary}

\theoremstyle{definition}
\newtheorem{definition}[theorem]{Definition}
\newtheorem{assumption}[theorem]{Assumption}

\theoremstyle{remark}
\newtheorem{remark}[theorem]{Remark}

\title{Bifurcation Bracketing via Asymptotic Stability Score}
\author{}
\date{\today}

\begin{document}
\maketitle


\section{Defintions}

\begin{definition}[System]
Our system is a nonlinear discrete dynamical system modelling predator-prey interactions with cyclic migration. For fixed Lotka–Volterra parameters
\[
\theta=(a,b,c,d,e,f)
\]and conflict parameter \(\alpha\), one iteration consists of the Lotka–Volterra update \(U_\theta\), normalization \(\mathcal N\), the conflict interaction \(C_\alpha\),
and inverse normalization \(\mathcal N^{-1}\) \citep{Albeverio01042008}. We therefore write the complete one-step transformation as
\[
F_\alpha
=
\mathcal N^{-1}\circ C_\alpha\circ\mathcal N\circ U_\theta.
\]For notational simplicity, the full transformation is denoted by \(F_\alpha\).
We define the state
\[
x=(\mathbf P,\mathbf R)
=(P_1,P_2,R_1,R_2)\in X\subseteq\mathbb R^4,
\]where \(\mathbf P=(P_1,P_2)\) and \(\mathbf R=(R_1,R_2)\) represent the two populations across the two patches \citep{Albeverio01042008}. The dynamics are then
\[
x_{t+1}=F_\alpha(x_t).
\]
During each simulation, both \(\theta\) and \(\alpha\) are held fixed. The parameter \(\alpha\) is varied only between simulations for the
purpose of bifurcation analysis,
giving the one-parameter family \(\{F_\alpha\}_{\alpha\in[-1,1]}\).
\end{definition}

\begin{definition}[Topological equivalence]
A dynamical system is called topologically equivalnt to another if there is a homeomorphism
\(h:\mathbb R^n\to\mathbb R^n\) mapping orbits of the first system to orbits of the second system,
preserving the direction of time \citep{Kuznetsov2023}.
\end{definition}

\begin{definition}[Conjugacy of maps]
Two maps \(f\) and \(g\) are called conjugate if there is a homeomorphism \(h:\mathbb R^n\to\mathbb R^n\) such that
\(h\circ f = g\circ h\) \citep{Kuznetsov2023}.
\end{definition}

\begin{definition}[Structural stability]
A system is structurally stable if there is a \(C^1\) neighborhood
\(\mathcal U\) of \(F\) such that every \(G \in \mathcal U\) is topologically conjugate to \(F\) \citep{Kuznetsov2023}.
\end{definition}

\begin{definition}[Bifurcation]
Kuznetsov (2023) defines a bifurcation as "The appearance of a topologically nonequivalent phase
portrait under variation of parameters" \citep{Kuznetsov2023}.
\end{definition}

\begin{definition}[stability score]
Let \((X, d)\) be a metric space where \(X\) is our system's state space and \(d\) is the \(L_1\) distance function where \(d(x, y) = \| x - y \|_1\).
We define the stability score as 
\[
S(\alpha) = \lim_{T \to \infty} \frac{1}{T} \sum_{t=0}^{T} d(x_{t+1}, x_t),
\]
where \(x_t = (p_t, r_t) \in \mathbb R^2\) is the state of the system at time \(t\).
\end{definition}

\section{Main Result}

\begin{theorem}
For \(F_\alpha(x) := F(x; \theta, \alpha)\) where \(\theta\) is fixed and exists \(\alpha, \alpha' \in [-1, 1]\) such that
\(\alpha < \alpha'\)
\(S(\alpha)\cdot S(\alpha') = 0\) and \(S(\alpha) + S(\alpha') > 0\) then exists bifurcation parameter \(\alpha^* \in [\alpha, \alpha']\).
\end{theorem}

The proof proceeds in three stages. First lemma 8 and lemma 9, show respectively that a fixed-point convergence implies that \(S(\alpha) = 0\)
and a non-trivial periodic attractor implies that \(S(\alpha) > 0\).
Then, lemma 10 shows that the endpoint attractors are not topologically conjugate. 
Finally, we combine these results to obtain the proof of theorem 7.

\begin{lemma}
A fixed-point convergence implies that S = 0.
\end{lemma}

\begin{proof}
Suppose that \(x_t \to x^*\) as \(t \to \infty\). Then \(x_{t+1} \to x^*\) as \(t \to \infty\). Therefore,
\[
d(x_{t+1}, x_t) \leq d(x_{t+1}, x^*) + d(x^*, x_t) \to 0
\]
as \(t \to \infty\). Therefore,
\[
\lim_{T \to \infty} \frac{1}{T} \sum_{t=0}^{T} d(x_{t+1}, x_t) = 0.
\]
Hence, \(S(\alpha) = 0\).
\end{proof}

\begin{lemma}
A non-trivial periodic attractor implies that S > 0.
\end{lemma}

\begin{proof}
Suppose the trajectory converges to a periodic orbit 
\[
\mathcal O = \{z_0, z_1, \ldots, z_{p-1}\}
\] for some \(p > 1\),
with \(F_{\alpha}(z_i) = z_{i+1}\) and \(F_{\alpha}(z_{p-1}) = z_0\) for
\(i = 0, 1, \ldots, p-1\).
Let \(T = qp + r\) such that \(0 \leq r < p\). Then,
\[
\frac{1}{T}\sum_{t=0}^{T} d(x_{t+1}, x_t) = \frac{1}{qp + r} 
(q\sum_{i=0}^{p-1} d(z_i, z_{i+1}) + \sum_{t=0}^{r-1} d(z_t, z_{t+1}))
\]
Given as \(T \to \infty\) then \(q \to \infty\) then,
\[
\lim_{T \to \infty} \frac{1}{T} \sum_{t=0}^{T} d(x_{t+1}, x_t) = 
\lim_{q \to \infty} \frac{1}{p + r/q} 
\sum_{i=0}^{p-1} d(z_i, z_{i+1}) + \frac{1}{qp + r}\sum_{t=0}^{r-1} d(z_t, z_{t+1}).
\]
Therefore,
\[
\lim_{T \to \infty} \frac{1}{T} \sum_{t=0}^{T} d(x_{t+1}, x_t) = 
\frac{1}{p}\sum_{i=0}^{p-1} d(z_i, z_{i+1}).
\]
Hence, given \(z_{i+1} \neq z_i\) then for all \(i = 0, 1, \ldots, p-1\) then \(S(\alpha) > 0\).

\end{proof}


\begin{lemma}
The endpoint attractors are not topologically conjugate.
\end{lemma}

\begin{proof}

Suppose there exists\(\alpha, \alpha' \in [-1, 1]\) such that
\(S(\alpha) = 0\) and \(S(\alpha') > 0\). Then, by lemma 1 
\(F_\alpha(x)\) has a fixed-point attractor \(A = \{x^*\}\) and by lemma 2
\(F_{\alpha'}(x)\) has a periodic attractor \(A' = \{z_0, z_1, \ldots, z_{p-1}\}\).

Assume there exists a conjugate homeomorphism 
\(h: A \to A'\) such that
\(h \circ F_\alpha = F_{\alpha'} \circ h\). Given \(F_\alpha(x^*) = x^*\) 
then \(h(F_\alpha(x^*)) = F_{\alpha'}(h(x^*))\) which results in
\(F_{\alpha'}(h(x^*)) = h(x^*)\).
Which implies \(h(x^*)\) is a fixed-point of \(F_{\alpha'}\) which contradicts the assumption that \(S(\alpha') > 0\).
Therefore, the endpoint attractors are not topologically conjugate.
\end{proof}

We can now combine the three lemmas to obtain the main result.

\begin{proof}
By lemma 8 \(S(\alpha) = 0\) implies the relevant attractor of \(F_\alpha\) 
is a fixed-point. By lemma 9 \(S(\alpha') > 0\) implies the relevant attractor of 
\(F_{\alpha'}\) is a non-trivial periodic orbit or vice-versa. 
By lemma 10, these two attractor dynamics are not topologically conjugate.

Suppose for contradiction no bifurcation parameter \(\alpha^* \in [\alpha, \alpha']\)
exists. Then the relevant attractor remains stable, so its topological type cannot
change continuously from a structurally stable regime to another without passing
through a parameter at which structural stability fails, hence
\(F_\alpha \sim F_{\alpha'}\) which contradicts lemma 10. 
Therefore, a bifurcation parameter \(\alpha^* \in [\alpha, \alpha']\) exists.
\end{proof}

\clearpage 
\bibliographystyle{apalike}
\bibliography{refs}

\end{document}
